\section{Markov chains}
\label{m:sec:chains}

A stochastic process is \emph{Markov} if its future depends on its past only
through its present. That single property is what makes every model in this thesis
tractable, and this section sets out the classes of Markov process the thesis uses.
Standard treatments are Norris \cite{Norris1997}, Karlin and Taylor
\cite{KarlinTaylor1975}, and, closest to the applications here, Allen
\cite{Allen2010}.

\subsection{Discrete-time Markov chains}
\label{m:sec:dtmc}

Let $\{X_n : n=0,1,2,\dots\}$ take values in a countable state space $\mathcal S$.
The process is a \emph{discrete-time Markov chain} if
\begin{equation}
\label{m:eq:dtmarkov}
\pr{X_{n+1}=j \mid X_n=i,\ X_{n-1},\dots,X_0} = \pr{X_{n+1}=j\mid X_n=i}
=: P_{ij},
\end{equation}
the right-hand side being independent of $n$ for a \emph{time-homogeneous} chain.
The matrix $P = (P_{ij})$ is stochastic: its entries are non-negative and each row
sums to one. Multi-step laws follow by decomposing a journey at an intermediate
state, which gives the Chapman--Kolmogorov identity
\begin{equation}
\label{m:eq:chapmankolmogorov}
\pr{X_{n+m}=j\mid X_0=i}
= \sum_{k\in\mathcal S}\pr{X_n=k\mid X_0=i}\,\pr{X_m=j\mid X_0=k},
\end{equation}
and hence $\pr{X_n=j\mid X_0=i}=(P^n)_{ij}$: composing steps is multiplying
matrices. In principle everything about the chain is therefore contained in $P$;
in practice the interest lies in extracting it, and the rest of this chapter is
about how. A state $i$ with $P_{ii}=1$ is \emph{absorbing}: a chain that reaches
it never leaves.

\subsubsection{Simple random walks}
\label{m:sec:randomwalk}

The simplest non-trivial chain moves one step at a time on the integers. Fix
$q\in(0,1)$ and let
\begin{equation}
\label{m:eq:rwrates}
P_{i,i+1} = q,
\qquad
P_{i,i-1} = 1-q,
\end{equation}
so that $X_n = X_0 + \sum_{k=1}^{n}\xi_k$ with the $\xi_k$ independent and equal to
$+1$ with probability $q$ and $-1$ otherwise. The walk is \emph{symmetric} when
$q=\tfrac12$. Its mean and variance are immediate from the independence of the
increments:
\begin{equation}
\label{m:eq:rwmoments}
\ex{X_n} = X_0 + n(2q-1),
\qquad
\var{X_n} = 4nq(1-q).
\end{equation}
The drift $2q-1$ decides the long-run behaviour, and the diffusive spread
$\sqrt{n}$ is what the mean cannot see.

Here the whole law is available too. If $R$ of the $n$
increments are $+1$ then $X_n = X_0+2R-n$, and $R$ is binomial with parameters
$n$ and $q$, so
\begin{equation}
\label{m:eq:rwlaw}
\pr{X_n = X_0+2r-n} = \binom{n}{r}q^{\,r}(1-q)^{\,n-r},
\qquad r=0,1,\dots,n ,
\end{equation}
supported on those lattice points with the parity of $n$
(\cref{m:fig:rwlaw}). The walk surrenders its $n$-step law so cheaply because its
increments are independent of the state it has reached. Once the step distribution
depends on the current count --- as it does for every branching process below ---
the $n$-step law is no longer a convolution and must be reached by the generating
functions of \cref{m:sec:pgfmethod}.

\begin{figure}[ht]
\centering
\includegraphics[width=0.74\textwidth]{rw_transition.pdf}
\caption{The exact displacement law \cref{m:eq:rwlaw} at $n=40$, for
$q=0.4,\,0.5,\,0.6$; the dotted vertical lines mark the means $(2q-1)n$ of
\cref{m:eq:rwmoments}. Changing the drift slides the whole distribution along the
lattice and barely deforms it, because the width $\sqrt{4nq(1-q)}$ varies little
over this range of $q$.}
\label{m:fig:rwlaw}
\end{figure}

The random walk is not a peripheral example here. Every continuous-time chain in
\cref{m:sec:ctmc} decomposes into a holding time and a jump, and the sequence of
jumps is itself a discrete-time chain --- the \emph{embedded jump chain}. For the
linear birth--death process of \cref{m:sec:bd}, whose rates out of state $i$ are
$\lambda i$ and $\mu i$, the state-dependence cancels in the ratio and the embedded
chain is exactly \cref{m:eq:rwrates} with
\begin{equation}
\label{m:eq:embedded}
q = \frac{\lambda}{\lambda+\mu} .
\end{equation}
Absorption at $0$ is then the walk's ruin problem, and hitting probabilities for
the continuous-time process can be computed on the walk alone.

Two absorbing barriers are usually imposed. With absorbing states at $0$ and $M$,
writing $h_i$ for the probability of reaching $M$ before $0$ from state $i$, the
one-step decomposition $h_i = qh_{i+1}+(1-q)h_{i-1}$ with $h_0=0$, $h_M=1$ has the
classical gambler's-ruin solution
\begin{equation}
\label{m:eq:gamblersruin}
h_i =
\begin{cases}
\dfrac{1-\theta^{\,i}}{1-\theta^{\,M}},
& \theta := \dfrac{1-q}{q}\neq1,\\[12pt]
\dfrac{i}{M}, & q=\tfrac12 .
\end{cases}
\end{equation}
This is the smallest instance of the first-step principle of
\cref{m:sec:firststep}, and \cref{m:fig:randomwalk} shows both halves of the
picture.

\begin{figure}[ht]
\centering
\includegraphics[width=0.98\textwidth]{random_walk.pdf}
\caption{The simple random walk of \cref{m:eq:rwrates}.
(a) Seven realisations at $q=0.45$ (grey), against the drift line
$n(2q-1)$ and the envelope $\pm\sqrt{4nq(1-q)}$ of one standard deviation
(red), from \cref{m:eq:rwmoments}.
(b) The gambler's-ruin probability \cref{m:eq:gamblersruin} of reaching $M=20$
before $0$, against the starting state $i$, for
$\theta=(1-q)/q \in\{0.7,\,0.85,\,1,\,1.2,\,1.5\}$. The symmetric case $\theta=1$
is the straight line $i/M$; $\theta<1$ bows the curve up and $\theta>1$ bows it
down, and the effect is strongly nonlinear in $\theta$.}
\label{m:fig:randomwalk}
\end{figure}

\subsubsection{Galton--Watson branching processes}
\label{m:sec:gw}

The Galton--Watson process is a discrete-time chain on the non-negative integers in
which every member of a population independently leaves offspring, or does not, at
each generational increment. It is the oldest model of stochastic branching and the
central object of this thesis; the standard references are Harris \cite{Harris1963}
and Athreya and Ney \cite{Athreya1972}. What follows is the definition and the two
or three identities the rest of this chapter uses. The process itself --- its
extinction dynamics, its critical endpoint, and the constant governing its
conditional mean --- is the subject of \ChA.

Suppose each individual produces, independently of every other, a number of
offspring distributed as the \rv{} $L$, with $\pr{L=k}=p_k$ and mean
$\mu=\ex{L}$. Let $Z_n$ be the size of generation $n$, started from a single
progenitor, $Z_0=1$, so that
\begin{equation}
\label{m:eq:gwrecursion}
Z_{n+1} = \sum_{i=1}^{Z_n}L_i ,
\end{equation}
the $L_i$ being independent copies of $L$. Conditioning on the previous generation
and using the tower property gives $\ex{Z_{n+1}} = \mu\ex{Z_n}$, hence
$\ex{Z_n}=\mu^{\,n}$, and the three regimes $\mu<1$, $\mu=1$, $\mu>1$ are called
\emph{subcritical}, \emph{critical} and \emph{supercritical}.

Because $Z_{n+1}$ is a random sum of independent copies of $L$, the composition
rule for generating functions (\cref{m:sec:pgfmethod}) turns
\cref{m:eq:gwrecursion} into a functional iteration. Writing
$\phi(z)=\ex{z^L}$ for the offspring generating function and $\phi_n$ for that of
$Z_n$,
\begin{equation}
\label{m:eq:gwcomposition}
\phi_{n+1} = \phi_n\circ\phi ,
\end{equation}
so $\phi_n$ is the $n$-fold composition of $\phi$ with itself. That single fact is
the source of both the tractability and the difficulty of the model: everything
about $Z_n$ is encoded in an iterated map, and iterating a map is easy to do and
hard to solve in closed form.

\paragraph{The binary case.}
Throughout this thesis we specialise to the binary, or \emph{simple}, case, in
which an individual either divides into two or dies:
\begin{equation}
\label{m:eq:binarylaw}
p_2 = p \quad\text{(division)},
\qquad
p_0 = 1-p \quad\text{(death)},
\qquad
p_i = 0 \ \text{ otherwise},
\end{equation}
so that $\mu=2p$, $\ex{Z_n}=(2p)^n$, and the offspring generating function is
\begin{equation}
\label{m:eq:pgfbinary}
\phi(z) = pz^2+(1-p) .
\end{equation}
Criticality is at $p=\tfrac12$. Writing $u_n = \pr{Z_n=0}$ for the probability of
extinction by generation $n$ and setting $z=0$ in \cref{m:eq:gwcomposition} gives
$u_{n+1}=\phi(u_n)$, and it is more convenient to work with the survival
probability $S_n = 1-u_n = \pr{Z_n>0}$, which satisfies
\begin{equation}
\label{m:eq:SnIter}
S_{n+1} = 2pS_n - pS_n^2,
\qquad S_0 = 1 .
\end{equation}

That is as far as the definition needs to go here. \Cref{m:eq:SnIter} is a quadratic
recursion, and almost everything of interest about the binary process turns out to
be a statement about its orbit: whether the orbit reaches zero, how fast it does
so, and what constant survives in front of the decay when it does. Those questions
--- extinction and its phase transition, the infinite expected lifetime at
criticality, and the amplitude $A(p)$ that governs the conditional mean --- are
taken up together in \ChA. What the present chapter needs of the process is the
definition above, the mean $\ex{Z_n}=(2p)^n$, the recursion \cref{m:eq:SnIter}, and
later a single formula for the conditional mean, given in \cref{m:sec:condmeans}
for comparison with its continuous-time counterpart.

\subsection{Continuous-time Markov chains}
\label{m:sec:ctmc}

Let $\{X_t : t\ge0\}$ take values in a countable state space $\mathcal S$, which
throughout will be a set of population counts or of tuples of counts. The process
is a \emph{continuous-time Markov chain} if, for all $s,t\ge0$,
\begin{equation}
\label{m:eq:markov}
\pr{X_{t+s}=j \mid X_t=i,\ \{X_u\}_{u<t}} = \pr{X_{t+s}=j\mid X_t=i},
\end{equation}
and \emph{time homogeneous} if the right-hand side depends on $s$ but not $t$, in
which case one writes $p_{ij}(s)$. Such a chain is specified by its
\emph{generator} $Q=(q_{ij})$, whose off-diagonal entries are transition rates,
\begin{equation}
\label{m:eq:generator}
q_{ij} = \lim_{\Delta t\to0}\frac{p_{ij}(\Delta t)}{\Delta t}\quad(j\neq i),
\qquad
q_{ii} = -\sum_{j\neq i}q_{ij} =: -q_i ,
\end{equation}
so that every row of $Q$ sums to zero. Equivalently, and this is the form used to
define models below,
\begin{equation}
\label{m:eq:infinitesimal}
\pr{X_{t+\Delta t}=j\mid X_t=i} = q_{ij}\,\Delta t + o(\Delta t)
\qquad(j\neq i),
\end{equation}
with complementary probability $1-q_i\Delta t+o(\Delta t)$ of no transition at all.
A state with $q_i=0$ is absorbing.

\subsubsection{Exponential clocks}
\label{m:sec:exp}

Everything in continuous time rests on one distribution, and it is worth seeing
why it is forced rather than chosen. A \rv{} $T$ is exponentially distributed with
rate $\theta>0$, written $T\sim\mathrm{Exp}(\theta)$, if
\begin{equation}
\label{m:eq:expdef}
f_T(t) = \theta\ee^{-\theta t},
\qquad
F_T(t) = 1-\ee^{-\theta t}
\qquad(t\ge0),
\end{equation}
with mean $1/\theta$ and variance $1/\theta^2$. The property that matters is
memorylessness: for $t_1,t_2>0$,
\begin{equation}
\label{m:eq:memoryless}
\pr{T>t_1+t_2\mid T>t_1}
= \frac{\ee^{-\theta(t_1+t_2)}}{\ee^{-\theta t_1}}
= \ee^{-\theta t_2}
= \pr{T>t_2},
\end{equation}
so a clock that has already run for time $t_1$ is statistically indistinguishable
from a fresh one. The exponential law is the only continuous distribution with this
property, which is exactly why a process assembled from constant-rate events
satisfies \cref{m:eq:markov}: with any other waiting-time law the elapsed time would
be informative and the process would not be Markov.

When several events compete, two facts organise everything that follows.

\begin{proposition}[Competing clocks]
\label{m:prop:race}
Let $T_1,\dots,T_m$ be independent with $T_i\sim\mathrm{Exp}(\theta_i)$, let
$T=\min_iT_i$, and write $\Theta=\sum_i\theta_i$. Then
\begin{equation}
\label{m:eq:race}
T\sim\mathrm{Exp}(\Theta),
\qquad
\pr{T_k=T} = \frac{\theta_k}{\Theta},
\end{equation}
and the identity of the winning clock is independent of $T$.
\end{proposition}

\begin{proof}
For the first claim,
$\pr{T>t} = \prod_i\pr{T_i>t} = \prod_i\ee^{-\theta_it} = \ee^{-\Theta t}$.
For the second, condition on the value of $T_k$ and use independence:
\begin{equation*}
\pr{T_k=T}
= \int_0^\infty\theta_k\ee^{-\theta_ku}\prod_{i\neq k}\ee^{-\theta_iu}\,\dd u
= \theta_k\int_0^\infty\ee^{-\Theta u}\,\dd u
= \frac{\theta_k}{\Theta} .
\end{equation*}
Repeating the calculation with the integral truncated at $t$ gives
$\pr{T_k=T,\ T>t} = (\theta_k/\Theta)\ee^{-\Theta t} = \pr{T_k=T}\pr{T>t}$, which is
the asserted independence.
\end{proof}

\Cref{m:prop:race} is what connects \cref{m:eq:infinitesimal} to a trajectory: the
chain leaves state $i$ after a holding time $\mathrm{Exp}(q_i)$ and then jumps to
$j$ with probability $q_{ij}/q_i$, independently of how long it waited. It also
says that a rate becomes a probability as soon as one asks \emph{which} of several
things happened rather than when.

\begin{remark}[Direct simulation]
\label{m:rem:gillespie}
\Cref{m:prop:race} is an algorithm as well as a theorem. Given the current state,
compute the rate $\theta_i$ of every possible event and their sum $\Theta$; draw
$T\sim\mathrm{Exp}(\Theta)$ by inverting \cref{m:eq:expdef} on a uniform variate;
draw the event $k$ with probability $\theta_k/\Theta$; advance the clock by $T$,
apply event $k$, and repeat. Because the two draws are independent and the
exponential law is memoryless, the resulting trajectory has exactly the law of the
process. This is Gillespie's direct method \cite{Gillespie1977}; it is used in this
thesis to check analytical results, never to replace them.
\end{remark}

\subsubsection{The Poisson process}
\label{m:sec:poisson}

The simplest continuous-time chain counts events that arrive at a constant rate and
never depart. Let $N_t$ take values in $\{0,1,2,\dots\}$ with the single non-zero
rate
\begin{equation}
\label{m:eq:poissonrates}
q_{n,n+1} = \alpha
\qquad(n\ge0),
\end{equation}
so $q_n=\alpha$ for every $n$. By \cref{m:prop:race} the holding times are
independent $\mathrm{Exp}(\alpha)$ variables, and $N_t$ counts how many have
elapsed by time $t$: the process is the renewal count of a sequence of independent
exponential waits. Its law follows from the master equation
\cref{m:eq:master} below, or directly by induction on $n$:
\begin{equation}
\label{m:eq:poissonlaw}
\pr{N_t = n} = \ee^{-\alpha t}\frac{(\alpha t)^n}{n!},
\qquad
\ex{N_t} = \var{N_t} = \alpha t .
\end{equation}
The equality of mean and variance is characteristic, and is the baseline against
which the over-dispersion of branching processes should be read.
\Cref{m:fig:poisson} shows a realisation and the law it is drawn from.

\begin{figure}[ht]
\centering
\includegraphics[width=0.98\textwidth]{poisson_process.pdf}
\caption{The Poisson process of \cref{m:eq:poissonrates}, at rate $\alpha=1.5$.
(a) One realisation of the counting process $N_t$, with the independent
$\mathrm{Exp}(\alpha)$ holding times of \cref{m:prop:race} marked in alternating
shades along the axis, and the mean $\alpha t$ dashed.
(b) The law \cref{m:eq:poissonlaw} at $t=1,\,3,\,6$. Mean and variance are both
$\alpha t$, so the distribution broadens as it advances and its relative width
$1/\sqrt{\alpha t}$ shrinks.}
\label{m:fig:poisson}
\end{figure}

Two properties are used later. Counts over disjoint time intervals are
independent, and the count over an interval depends only on its length --- the
process has stationary independent increments, which is the continuous-time
statement of memorylessness. If each arrival is independently retained with
probability $\varpi$, the retained arrivals form a Poisson process of rate
$\varpi\alpha$: \emph{thinning} preserves the class.

\subsubsection{Birth--death processes}
\label{m:sec:bd}

Let $X_t$ count individuals, each of which independently gives birth at rate
$\lambda$ and dies at rate $\mu$. The non-zero rates out of state $n$ are
\begin{equation}
\label{m:eq:bdrates}
q_{n,n+1} = \lambda n,
\qquad
q_{n,n-1} = \mu n ,
\end{equation}
so $q_n = (\lambda+\mu)n$ and $n=0$ is absorbing. This is the \emph{linear}
birth--death process, the continuous-time counterpart of the Galton--Watson process
of \cref{m:sec:gw}, and it is the running example of \cref{m:sec:ctmethods}: nearly
every method in this thesis is first demonstrated on it. Its embedded jump chain is
the random walk \cref{m:eq:rwrates} with $q=\lambda/(\lambda+\mu)$, so questions
about \emph{which} states are visited reduce to \cref{m:sec:randomwalk} while
questions about \emph{when} require the rates themselves.

\Cref{m:fig:bdpaths} shows the three regimes it can occupy.

\begin{figure}[ht]
\centering
\includegraphics[width=0.98\textwidth]{birth_death_paths.pdf}
\caption{Nine realisations of the linear birth--death process \cref{m:eq:bdrates}
from $X_0=5$, with $\mu=1$ throughout and the mean $N\ee^{(\lambda-\mu)t}$ of
\cref{m:eq:bdmean} in red; dots on the axis mark trajectories that reached the
absorbing state. The middle panel is the one to dwell on. At criticality the mean
is exactly constant, and yet several realisations are already extinct while others
have grown fourfold --- the mean is the same number in both cases, and describes
neither. Note also that extinction is not confined
to the subcritical panel: a supercritical process still dies out with probability
$(\mu/\lambda)^{N}$, visible here as the paths that fail early on the right.}
\label{m:fig:bdpaths}
\end{figure}

The general birth--death process allows state-dependent rates $\lambda_n$ and
$\mu_n$; the linear case $\lambda_n=\lambda n$, $\mu_n=\mu n$ is the one in which
individuals act independently, and hence the one for which generating functions
close.

\subsubsection{Birth--death--catastrophe processes}
\label{m:sec:bdc}

The processes at the centre of this thesis add a third channel. Alongside the
gradual loss of individuals to death, a compartment may suffer a
\emph{catastrophe}: an abrupt event that ends the entire population at once.

\begin{definition}[Birth--death--catastrophe process]
\label{m:def:bdc}
Let $X_t$ take values in $\{0,1,2,\ldots\}\cup\{H\}$ and count individuals in a
compartment while it remains intact. Each individual independently gives birth at
rate $\lambda$, dies at rate $\mu$, and triggers a catastrophe at rate $\delta$.
Once any individual triggers that event, the whole compartment enters the
post-catastrophe state $H$. The non-zero transition rates out of $n\ge1$ are
\begin{equation}
\label{m:eq:bdcrates}
q_{n,n+1} = \lambda n,
\qquad
q_{n,n-1} = \mu n,
\qquad
q_{n,H} = \delta n ,
\end{equation}
and both $0$ and $H$ are absorbing.
\end{definition}

The two absorbing states record different terminal fates. State $0$ is reached by
the gradual deaths of all individuals; a separate state $H$ preserves the
distinction between that internal extinction and catastrophe, including any load
released at rupture. (In the infection application, $H$ is host-cell rupture.)
If $\tau$ is the catastrophe time, the load released is the pre-catastrophe count
$X_{\tau-}$; \fwd{bdc}{the chapter on the birth--death--catastrophe process}
records this load as a second process. Because the
total catastrophe rate in state $n$ is $\delta n$, a fuller compartment is more
likely to rupture: the hazard is state-dependent. Up to the first catastrophe,
family trees and their individual trigger clocks evolve independently, so
\emph{no-catastrophe} probabilities for several founders factorise into
single-founder contributions: from $k$ founders the probability that none has
triggered by time $t$ is the $k$th power of its single-founder value. Events
that reach past the first trigger do not factorise in this way, since one
trigger stops every lineage at once.
\Cref{m:fig:ratediagram} sets the linear birth--death and
birth--death--catastrophe processes side by side.

\begin{figure}[ht]
\centering
\begin{tikzpicture}[
  >={Stealth[length=1.9mm]},
  st/.style ={circle, draw, thick, minimum size=7.5mm, inner sep=0pt, font=\small},
  ab/.style ={circle, draw, thick, double, double distance=0.7pt,
              minimum size=7.5mm, inner sep=0pt, font=\small},
  lb/.style ={font=\footnotesize, inner sep=1.6pt},
  rowlab/.style={font=\small\itshape, anchor=west},
  every path/.style={thick}
]
% ---------------- linear birth--death -------------------------------------
\node[rowlab] at (-1.9,1.15) {birth--death};
\node[ab] (a0) at (0,0)   {$0$};
\node[st] (a1) at (2.1,0) {$1$};
\node[st] (a2) at (4.2,0) {$2$};
\node[st] (a3) at (6.3,0) {$3$};
\node     (a4) at (7.9,0) {$\cdots$};
\draw[->] (a1) to[bend left=38]  node[above,lb]{$\lambda$}  (a2);
\draw[->] (a2) to[bend left=38]  node[above,lb]{$2\lambda$} (a3);
\draw[->] (a3) to[bend left=38]  node[above,lb]{$3\lambda$} (a4);
\draw[->] (a1) to[bend left=38]  node[below,lb]{$\mu$}      (a0);
\draw[->] (a2) to[bend left=38]  node[below,lb]{$2\mu$}     (a1);
\draw[->] (a3) to[bend left=38]  node[below,lb]{$3\mu$}     (a2);

% ---------------- birth--death--catastrophe -------------------------------
\begin{scope}[yshift=-3.5cm]
\node[rowlab] at (-1.9,1.15) {birth--death--catastrophe};
\node[ab] (b0) at (0,0)   {$0$};
\node[st] (b1) at (2.1,0) {$1$};
\node[st] (b2) at (4.2,0) {$2$};
\node[st] (b3) at (6.3,0) {$3$};
\node     (b4) at (7.9,0) {$\cdots$};
\draw[->] (b1) to[bend left=38]  node[above,lb]{$\lambda$}  (b2);
\draw[->] (b2) to[bend left=38]  node[above,lb]{$2\lambda$} (b3);
\draw[->] (b3) to[bend left=38]  node[above,lb]{$3\lambda$} (b4);
\draw[->] (b1) to[bend left=38]  node[below,lb]{$\mu$}      (b0);
\draw[->] (b2) to[bend left=38]  node[below,lb]{$2\mu$}     (b1);
\draw[->] (b3) to[bend left=38]  node[below,lb]{$3\mu$}     (b2);
\node[ab] (bd) at (4.2,-2.35) {$H$};
\draw[->,densely dashed] (b1) to[bend right=12]
      node[pos=0.68,left=1.5pt,lb]{$\delta$} (bd);
\draw[->,densely dashed] (b2) --
      node[pos=0.68,right=1.5pt,lb]{$2\delta$} (bd);
\draw[->,densely dashed] (b3) to[bend left=12]
      node[pos=0.68,right=1.5pt,lb]{$3\delta$} (bd);
\end{scope}
\end{tikzpicture}
\caption{The linear birth--death process \cref{m:eq:bdrates}, above, and the
birth--death--catastrophe process of \cref{m:def:bdc}, below. Double circles are
absorbing. In both rows the only route into state $0$ is a death out of state
$1$, so ordinary extinction must pass through the smallest population;
catastrophe (dashed) instead reaches the separate state $H$ from any $n\ge1$, at
total rate $\delta n$, so a fuller compartment is the more likely to rupture.}
\label{m:fig:ratediagram}
\end{figure}

\fwd{bdc}{The chapter on the birth--death--catastrophe process} develops the
process of \cref{m:def:bdc}: the probabilities of
rupture and of internal extinction, and the moments of the intracellular count
and of the load released at catastrophe. The remainder of the present chapter
supplies the general Markov and generating-function tools used there; it does not
derive the BDC-specific results themselves.

\begin{remark}[A state-dependent catastrophe is not a mean-field hazard]
\label{m:rem:notmeanfield}
It is tempting to replace a state-dependent catastrophe rate by a deterministic
hazard evaluated at the mean population, and it does not work. The point is
clearest in discrete time. Write $\chi\in(0,1)$ for the probability that a given
individual initiates catastrophe at a generation. It is the discrete counterpart
of the per-capita catastrophe rate $\delta$ of \cref{m:def:bdc}, whose total rate
out of census $n$ is $\delta n$ into the absorbing rupture state $H$; the two
keep separate letters throughout, one being a probability per generation and the
other a rate per unit time. Let
$\widetilde X_i$ be the branching population with catastrophe switched off, and
let
\begin{equation*}
\mathcal{E}_n = \sum_{i=0}^{n-1}\widetilde X_i
\end{equation*}
be the cumulative exposure through generation $n-1$ (the letter $H$ is reserved
for the catastrophe state itself). Conditioning on the entire catastrophe-free
genealogy, every individual present independently declines to trigger, so
\begin{equation}
\label{m:eq:norupture}
\pr{\text{no catastrophe before }n} = \ex{(1-\chi)^{\mathcal{E}_n}} .
\end{equation}
This identity is exact. Replacing $\mathcal{E}_n$ by $\ex{\mathcal{E}_n}$ is not:
the map $e\mapsto(1-\chi)^e$ is convex, so Jensen's inequality gives
$\ex{(1-\chi)^{\mathcal{E}_n}}\ge(1-\chi)^{\ex{\mathcal{E}_n}}$, and the
substitution systematically understates the chance of getting through: the
expectation in \cref{m:eq:norupture} is a nonlinear functional of the whole path,
and no deterministic hazard closes it.
Two further cautions belong with it. Non-catastrophe is not the same event as
population survival: the probability of being both unruptured and non-zero at
generation $n$ is $\ex{\ind\{\widetilde X_n>0\}(1-\chi)^{\mathcal{E}_n}}$, for
which Jensen supplies no bound at all. And once a trigger occurs every surviving
lineage is stopped together --- the compartment enters $H$ --- so the killed
lineages are coupled and cannot be treated as independent after the fact ---
although their family trees and trigger variables do factorise \emph{before} any
trigger, which is what makes \cref{m:eq:norupture} correct. Closing the model
properly therefore requires a marked or multivariate generating function, or the
rate description of \cref{m:def:bdc,m:eq:bdcrates}.
\end{remark}

\subsection{Time-inhomogeneous processes}
\label{m:sec:inhomogeneous}

All the chains above are time homogeneous: the rates in \cref{m:eq:infinitesimal}
depend on the state but not on the clock. Dropping that assumption gives a
\emph{time-inhomogeneous} chain, with generator $Q(t)$ and transition
probabilities depending on both endpoints,
\begin{equation}
\label{m:eq:inhomogtrans}
p_{ij}(s,t) = \pr{X_t=j\mid X_s=i}
\qquad(s\le t),
\end{equation}
in place of the single-argument $p_{ij}(t-s)$. The Chapman--Kolmogorov identity
survives in the form $P(s,t) = P(s,u)P(u,t)$ for $s\le u\le t$, and differentiating
it gives forward and backward equations
\begin{equation}
\label{m:eq:inhomogkolmogorov}
\pdd{}{t}P(s,t) = P(s,t)\,Q(t),
\qquad
\pdd{}{s}P(s,t) = -Q(s)\,P(s,t),
\end{equation}
but the solution is no longer the matrix exponential $\ee^{(t-s)Q}$ unless the
generators at different times commute. In general one is left with a
time-ordered product, and the methods of \cref{m:sec:ctmethods} apply only after
that complication has been handled.

Much survives nonetheless. The competing-clocks decomposition of
\cref{m:prop:race} still holds instantaneously, so direct simulation still applies
with a time-dependent total rate --- at the cost of thinning or of integrating the
rate to find the next event time. And a linear birth--death process with
time-varying rates $\lambda(t)$, $\mu(t)$ still has a generating function
satisfying a first-order partial differential equation, so the method of
characteristics of \cref{m:sec:moc} remains available; the characteristics simply
acquire time-dependent coefficients.

\subsubsection{The time-inhomogeneous linear birth--death process}
\label{m:sec:tibd}

The workhorse of the inhomogeneous setting is the process of \cref{m:sec:bd} with
rates that are prescribed functions of the clock: each individual gives birth at
rate $\lambda(t)$ and dies at rate $\mu(t)$, both taken locally integrable. The
master equation is \cref{m:eq:bdmaster} with time-dependent coefficients, and the
generating function accordingly satisfies
\begin{equation}
\label{m:eq:tibdpde}
\pddt{G} = \bigl(\lambda(t)z^2-(\lambda(t)+\mu(t))z+\mu(t)\bigr)\pdd{G}{z},
\qquad
G(z,0)=z^N ,
\end{equation}
which is \cref{m:eq:bdpde} with its quadratic coefficient set in motion. The
equation is still first order, so \cref{m:sec:moc} still applies; what has been
lost is that the coefficient no longer factorises once and for all, and with it
the closed solution \cref{m:eq:bdsolution}.

Two quantities survive intact whatever the rates do. The mean obeys the same
closed equation as before, $\ddt{}\ex{X_t}=(\lambda(t)-\mu(t))\ex{X_t}$, because
the rates are still linear in the count, so
\begin{equation}
\label{m:eq:meanodeinhom}
\ex{X_t} = N\exp\lrb{\int_0^t\bigl(\lambda(s)-\mu(s)\bigr)\dd s} ,
\end{equation}
and only the integrated net growth enters, not how it was distributed over the
interval. And the extinction probability $u(t)=\pr{X_t=0}$ from a single founder
satisfies the Riccati equation
\begin{equation}
\label{m:eq:riccatiext}
\dda{u}{t} = \mu(t)-\bigl(\lambda(t)+\mu(t)\bigr)u+\lambda(t)u^2,
\qquad u(0)=0 ,
\end{equation}
which is the backward equation of \cref{m:sec:backward} with time-dependent
coefficients: over the first small interval either the founder dies, contributing
$\mu(t)\Delta t$, or it divides, contributing $\lambda(t)\Delta t\,u(t)^2$ since
both daughter lineages must then die out in turn, or nothing happens. This is a
single scalar equation, integrated numerically without difficulty when it does
not solve in closed form, and the full generating function retains the
fractional-linear shape of \cref{m:eq:bdsolution} with its coefficients replaced
by solutions of \cref{m:eq:riccatiext}.

\subsubsection{The logistic speciation model}
\label{m:sec:logisticspeciation}

The standard application of that framework is a population whose per-capita birth
rate falls as the population fills the space available to it. In the macroevolutionary
reading of that idea the counted individuals are species and the births are
speciation events \cite{Allen2010}: let $X_t$ be the standing diversity, let each
species go extinct at constant rate $\mu$, and let it speciate at the
diversity-dependent rate
\begin{equation}
\label{m:eq:speciationrate}
\lambda(n) = \lambda_0\lrb{1-\frac{n}{K}},
\qquad \lambda_0>\mu ,
\end{equation}
when the standing diversity is $n$, with $K$ the niche capacity.

\Cref{m:eq:speciationrate} is not yet of the form treated above, and the
distinction is the whole point of the subsection: the rate depends on the
\emph{state}, not on the clock. The chain it defines is nonlinear, its moment
equations do not close, and none of the machinery of \cref{m:sec:ctmethods}
reaches it directly. The tractable surrogate replaces the random diversity inside
\cref{m:eq:speciationrate} by a deterministic mean-field trajectory $\bar N(t)$,
which turns the rate into a prescribed function of time and the chain into a
time-inhomogeneous \emph{linear} process, to which \cref{m:sec:tibd} does apply.
That replacement is a modelling step and not an identity: the mean is substituted
for a random census, and the resulting process is a comparison object rather than
an exact reduction of the nonlinear chain. Everything that follows is accordingly
a statement about the surrogate, not about the chain \cref{m:eq:speciationrate}
defines.

The mean-field diversity solves the logistic equation
\begin{equation}
\label{m:eq:logisticode}
\dda{\bar N}{t} = \bigl(\lambda(\bar N)-\mu\bigr)\bar N
= \gamma\,\bar N\lrb{1-\frac{\bar N}{N_\ast}},
\qquad
\gamma = \lambda_0-\mu,
\quad
N_\ast = K\lrb{1-\frac{\mu}{\lambda_0}},
\end{equation}
whose equilibrium $N_\ast$ is the diversity at which speciation balances
extinction. From $\bar N(0)=N_0$,
\begin{equation}
\label{m:eq:logisticsol}
\bar N(t) = \frac{N_\ast}{1+B\ee^{-\gamma t}},
\qquad
B = \frac{N_\ast}{N_0}-1 ,
\end{equation}
and the surrogate speciation rate $\lambda(t)=\lambda_0\bigl(1-\bar N(t)/K\bigr)$
is
\begin{equation}
\label{m:eq:lambdat}
\lambda(t) = \mu + \frac{\gamma B\ee^{-\gamma t}}{1+B\ee^{-\gamma t}},
\qquad
\lambda(0) = \lambda_0\lrb{1-\frac{N_0}{K}} ,
\end{equation}
relaxing monotonically from $\lambda(0)$ down to $\mu$. The integrated net growth
that \cref{m:eq:meanodeinhom} calls for is then available in closed form,
\begin{equation}
\label{m:eq:lambdaint}
\int_0^t\bigl(\lambda(s)-\mu\bigr)\dd s
= \ln\lrb{\frac{1+B}{1+B\ee^{-\gamma t}}} ,
\end{equation}
so that the mean diversity of the surrogate process is
\begin{equation}
\label{m:eq:logspecmean}
\ex{X_t} = N_0\,\frac{1+B}{1+B\ee^{-\gamma t}}
\;\xrightarrow[t\to\infty]{}\;
N_0(1+B) = N_\ast ,
\end{equation}
recovering the deterministic equilibrium, as a linear surrogate built on that
equilibrium must. Its extinction probability follows by integrating
\cref{m:eq:riccatiext} with the rate \cref{m:eq:lambdat}, and no longer in closed
form. \Cref{m:fig:logspecmean} plots \cref{m:eq:logspecmean} for three founding
diversities.

\begin{figure}[ht]
\centering
\includegraphics[width=0.78\textwidth]{logspec_mean.pdf}
\caption{Mean diversity \cref{m:eq:logspecmean} of the logistic speciation
surrogate for founding diversities $N_0=1,3,6$, with $\mu/\lambda_0=0.1$ and
$K=10$, so that $N_\ast=9$; the abscissa is the rescaled time $\gamma t$. Every
founding diversity relaxes monotonically to $N_\ast$, and the mean of the
surrogate coincides with the deterministic solution \cref{m:eq:logisticsol}
exactly, because the surrogate is linear and its mean therefore closes. The
smallest cohort spends a long early interval in near-exponential growth before the
diversity dependence binds --- and it is over precisely that interval that the
surrogate is least trustworthy, since it is there that a single realisation is
least like its mean.}
\label{m:fig:logspecmean}
\end{figure}

Two limits are worth reading off \cref{m:eq:lambdat}. At early times
$\lambda(t)\approx\lambda(0)>\mu$ and the surrogate is indistinguishable from a
homogeneous supercritical process; at late times $\lambda(t)\downarrow\mu$ and it
approaches criticality from above. The model therefore drifts, over its own
lifetime, into the regime whose discrete-generation counterpart causes the
difficulties documented in \cref{m:app:smallpops} --- a fixed carrying capacity
manufactures a near-critical process out of a comfortably supercritical one, and
does so without any parameter being changed by hand.

\subsection{Coupled ODE--CTMC systems}
\label{m:sec:coupled}

The final class leaves the purely stochastic setting. In a \emph{coupled
ODE--CTMC system} a continuous, deterministic state $\mathbf y(t)\in\mathbb R^d$
evolves alongside a discrete, stochastic state $X_t\in\mathcal S$, and the two
influence one another. Three couplings are possible, and they are genuinely
different.

\begin{enumerate}[label=(\roman*),leftmargin=*]
\item \emph{The differential equation depends on the chain.} The vector field is
      indexed by the discrete state,
      \begin{equation}
      \label{m:eq:coupledone}
      \dda{\mathbf y}{t} = F\bigl(\mathbf y(t),\,X_t\bigr),
      \end{equation}
      with $X_t$ an autonomous Markov chain. Between jumps the trajectory follows a
      fixed flow; at each jump the flow switches. The pair
      $(\mathbf y(t),X_t)$ is a Markov process on $\mathbb R^d\times\mathcal S$ even
      though neither component is one alone, and it is a piecewise-deterministic
      Markov process in the standard classification \cite{Davis1993}.
\item \emph{The chain depends on the differential equation.} The transition rates
      are functions of the continuous state,
      \begin{equation}
      \label{m:eq:coupledtwo}
      q_{ij} = q_{ij}\bigl(\mathbf y(t)\bigr),
      \end{equation}
      with $\mathbf y$ following a prescribed trajectory. This is a
      time-inhomogeneous chain in the sense of \cref{m:sec:inhomogeneous}, with the
      inhomogeneity generated by the flow rather than imposed directly, and the
      remarks there apply.
\item \emph{Both directions at once.} The vector field depends on the discrete
      state and the rates depend on the continuous one:
      \begin{equation}
      \label{m:eq:coupledthree}
      \dda{\mathbf y}{t} = F\bigl(\mathbf y(t),X_t\bigr),
      \qquad
      q_{ij} = q_{ij}\bigl(\mathbf y(t)\bigr) .
      \end{equation}
      Neither component can now be solved first and substituted into the other, and
      the joint process must be treated as a whole.
\end{enumerate}

Case (i) alone is already enough to break the methods of \cref{m:sec:ctmethods} as
stated, because the state space is no longer countable and the generator acquires a
transport term: for suitable test functions $f(\mathbf y,i)$,
\begin{equation}
\label{m:eq:coupledgenerator}
(\mathcal Lf)(\mathbf y,i)
= F(\mathbf y,i)\cdot\nabla_{\mathbf y}f(\mathbf y,i)
+ \sum_{j\neq i}q_{ij}(\mathbf y)\bigl\{f(\mathbf y,j)-f(\mathbf y,i)\bigr\} .
\end{equation}
The first term is the deterministic drift along the flow and the second is the
familiar jump part; \cref{m:eq:coupledgenerator} covers all three cases above,
with the $\mathbf y$-dependence of the rates dropped in case (i) and the
$i$-dependence of $F$ dropped in case (ii). Written this way the coupled system is
a recognisable object, and the corresponding forward equation is a system of
transport equations indexed by $\mathcal S$ rather than a system of ordinary
differential equations.

\Cref{m:fig:coupled} shows what the two-way case \cref{m:eq:coupledthree} looks like
in a trajectory. The system drawn there is generic --- it is chosen only to make the
mechanism visible, and is not a model used elsewhere in this thesis --- but it
displays the feature that distinguishes coupling (iii) from the two one-way cases.
Because the switching rate rises with $y$, the intervals spent in the second
discrete state are not scattered uniformly in time: they cluster where the
continuous variable has climbed, and each such interval then pulls $y$ back down.
Neither component can be solved first and substituted into the other, because each
is continually re-shaping the other's behaviour.

\begin{figure}[ht]
\centering
\includegraphics[width=0.98\textwidth]{coupled_ode_ctmc.pdf}
\caption{A trajectory of a coupled ODE--CTMC system of the two-way form
\cref{m:eq:coupledthree}. The continuous variable follows
$\dd y/\dd t = a(X_t)-by$, so it relaxes towards $a_1/b$ while $X_t=1$ and towards
$a_2/b$ while $X_t=2$ (dashed lines); the chain switches $1\to2$ at rate
$\kappa y$, which depends on the continuous state, and back at a constant rate.
Upper panel: $y(t)$, with the intervals spent in state $2$ shaded. Lower panel:
the discrete component $X_t$. The visits to state $2$ concentrate where $y$ is
large, which is the $y$-dependence of the rate made visible; the resulting descents
then remove the very condition that made them likely. The joint process
$(y(t),X_t)$ is Markov although neither component is, and its generator is
\cref{m:eq:coupledgenerator}. The parameters are illustrative only.}
\label{m:fig:coupled}
\end{figure}

The class is not introduced for its own sake. A ruptured compartment does not
simply vanish: at the catastrophe time $\tau$ of \cref{m:def:bdc} its contents
are released at once into the medium around it, so the continuous variable
receives an impulsive jump rather than a smooth contribution, and the size of
that jump is the pre-catastrophe count $X_{\tau-}$ --- the state of the chain at
a stopping time, not a parameter of the model. Both arrows of
\cref{m:eq:coupledthree} are then present, since a fuller medium may in turn
hasten the next rupture. The models built on that mechanism belong to
\fwd{rupture}{the chapter on rupture into a shared medium}; what the present
section supplies is the class they inhabit.

\subsubsection{Simulating a coupled system}
\label{m:sec:coupledsim}

The competing-clocks construction of \cref{m:prop:race} extends to
\cref{m:eq:coupledthree} with one qualification: the rates are no longer constant
between events, because $\mathbf y$ moves. Given the current state
$(\mathbf y,i)$, integrate the flow with the discrete component held at $i$ to
obtain $\mathbf y_s$ for $s\ge t$; the total jump rate
$q_i(\mathbf y_s)=\sum_{j\neq i}q_{ij}(\mathbf y_s)$ is then a known function of
$s$, and the next jump time $T$ has survival function
\begin{equation}
\label{m:eq:coupledholding}
\pr{T>s} = \exp\lrb{-\int_t^{s}q_i(\mathbf y_u)\,\dd u} ,
\end{equation}
exactly as in \cref{m:sec:inhomogeneous}; the destination is then drawn with
probability $q_{ij}(\mathbf y_T)/q_i(\mathbf y_T)$, the state is updated, and the
procedure repeats. Where $q_i$ admits a constant majorant the integral can be
inverted by rejection, and otherwise by quadrature. Analytically the position is
less comfortable. If the coupling is weak one closes the system by replacing
$\mathbf y$ by its mean-field trajectory --- which returns a time-inhomogeneous
chain in the sense of \cref{m:sec:inhomogeneous}, and the logistic surrogate of
\cref{m:sec:logisticspeciation} is exactly that closure --- or by replacing $X_t$
by its mean, which returns a closed differential equation. When the coupling is
strong neither closure is exact, and the working objects become joint
expectations satisfying hybrid moment equations that mix the jump generator with
the drift.
